\chapter{La línea recta}

\begin{center}
\begin{minipage}{.88\textwidth}
\small\sffamily
\textbf{Propósito.} Representar rectas mediante ecuaciones, interpretar sus
parámetros y resolver problemas de posición, intersección, ángulo y distancia.
\end{minipage}
\end{center}

\begin{properties}
Al terminar la unidad podrás:
\begin{itemize}
  \item calcular e interpretar la pendiente y el ángulo de inclinación;
  \item obtener la ecuación de una recta a partir de distintos datos;
  \item convertir una ecuación entre sus formas más utilizadas;
  \item decidir si dos rectas son paralelas, perpendiculares o secantes;
  \item calcular intersecciones, ángulos y distancias.
\end{itemize}
\end{properties}

\section{Pendiente y ángulo de inclinación}

\begin{definition}{Pendiente de una recta}{pendiente}
Sean $P_1(x_1,y_1)$ y $P_2(x_2,y_2)$ dos puntos distintos de una recta, con
$x_1\neq x_2$. La pendiente es la razón entre el cambio vertical y el cambio
horizontal:
\[
m=\frac{\Delta y}{\Delta x}.
\]
\end{definition}

\begin{formula}[title={Pendiente entre dos puntos}]
\[
\boxed{m=\frac{y_2-y_1}{x_2-x_1}},\qquad x_1\neq x_2.
\]
\end{formula}

\begin{intuition}
La pendiente indica cuánto cambia $y$ cuando $x$ aumenta una unidad. Si
$m=2$, por cada unidad hacia la derecha se suben dos; si $m=-\frac12$, por
cada dos unidades hacia la derecha se baja una.
\end{intuition}

\begin{properties}
\begin{center}
\begin{tabular}{ccl}
\toprule
Pendiente & Comportamiento & Interpretación\\
\midrule
$m>0$ & creciente & sube de izquierda a derecha\\
$m<0$ & decreciente & baja de izquierda a derecha\\
$m=0$ & horizontal & $y$ permanece constante\\
no definida & vertical & $x$ permanece constante\\
\bottomrule
\end{tabular}
\end{center}
\end{properties}

\begin{example}{Cálculo e interpretación de la pendiente}{calcular-pendiente}
Calcula la pendiente de la recta que pasa por $A(-2,1)$ y $B(4,7)$.
\begin{solution}
\textbf{Paso 1.} Identificamos $x_1=-2$, $y_1=1$, $x_2=4$ y $y_2=7$.

\textbf{Paso 2.} Calculamos los cambios:
\[
\Delta x=4-(-2)=6,\qquad \Delta y=7-1=6.
\]

\textbf{Paso 3.} Formamos la razón:
\[
m=\frac{\Delta y}{\Delta x}=\frac66=1.
\]

\textbf{Paso 4.} Interpretamos: al avanzar una unidad en $x$, la altura
aumenta una unidad. La recta es creciente.
\end{solution}
\end{example}

\begin{remark}
Al intercambiar los puntos cambian de signo tanto el numerador como el
denominador, de modo que la pendiente no cambia.
\end{remark}

\subsection{Ángulo de inclinación}

El ángulo de inclinación $\theta$ se mide en sentido antihorario desde el eje
$x$ positivo hasta la recta, con $0^\circ\leq\theta<180^\circ$.

\begin{formula}[title={Relación entre pendiente e inclinación}]
\[
\boxed{m=\tan\theta},
\qquad
\boxed{\theta=\arctan(m)}.
\]
Si $m<0$, se suma $180^\circ$ al valor principal de $\arctan(m)$ para obtener
un ángulo en el intervalo indicado.
\end{formula}

\begin{example}{Ángulo de inclinación}{inclinacion}
Determina el ángulo de inclinación de una recta cuya pendiente es
$m=-\sqrt3$.
\begin{solution}
\textbf{Paso 1.} Calculamos el valor principal:
\[
\arctan(-\sqrt3)=-60^\circ.
\]

\textbf{Paso 2.} Como el ángulo de inclinación debe estar entre $0^\circ$ y
$180^\circ$, sumamos $180^\circ$:
\[
\theta=-60^\circ+180^\circ=120^\circ.
\]
\end{solution}
\end{example}

\begin{exercise}{Pendiente e inclinación}{ej-pendiente}
\begin{enumerate}[label=\alph*)]
  \item Calcula la pendiente entre $A(1,3)$ y $B(5,11)$.
  \item Calcula la pendiente entre $C(-4,6)$ y $D(2,-3)$.
  \item Determina si la recta por $P(-2,5)$ y $Q(4,5)$ es creciente,
  decreciente, horizontal o vertical.
  \item ¿Qué ocurre con la fórmula si los puntos son $R(3,-1)$ y $S(3,8)$?
  \item Obtén el ángulo de inclinación para $m=1$, $m=\sqrt3$ y $m=-1$.
\end{enumerate}
\end{exercise}

\section{Ecuaciones de la recta}

Una ecuación de dos variables representa una recta cuando todos los puntos de
la recta la satisfacen y todo punto que la satisface pertenece a ella.

\subsection{Forma punto--pendiente}

Si una recta de pendiente $m$ pasa por $P_1(x_1,y_1)$, cualquier otro punto
$P(x,y)$ de la recta cumple
\[
m=\frac{y-y_1}{x-x_1}.
\]
Al multiplicar por $x-x_1$ se obtiene la forma siguiente.

\begin{formula}[title={Forma punto--pendiente}]
\[
\boxed{y-y_1=m(x-x_1)}.
\]
\end{formula}

\begin{example}{Recta con un punto y la pendiente}{punto-pendiente}
Obtén la ecuación de la recta de pendiente $m=-2$ que pasa por $P(3,5)$.
\begin{solution}
\textbf{Paso 1.} Sustituimos $m=-2$, $x_1=3$ y $y_1=5$:
\[
y-5=-2(x-3).
\]

\textbf{Paso 2.} Desarrollamos:
\[
y-5=-2x+6.
\]

\textbf{Paso 3.} Despejamos $y$:
\[
\boxed{y=-2x+11}.
\]

\textbf{Paso 4.} Comprobamos con $P(3,5)$:
$5=-2(3)+11=5$.
\end{solution}
\end{example}

\subsection{Forma pendiente--ordenada al origen}

\begin{formula}[title={Forma pendiente--ordenada}]
\[
\boxed{y=mx+b}.
\]
$m$ es la pendiente y $b$ es la ordenada al origen, es decir, la coordenada
$y$ del punto donde la recta corta el eje vertical.
\end{formula}

\begin{example}{Interpretación de $m$ y $b$}{pendiente-ordenada}
Analiza y grafica la recta $y=\frac12x-2$.
\begin{solution}
\textbf{Paso 1.} Identificamos $m=\frac12$ y $b=-2$.

\textbf{Paso 2.} Marcamos el corte con el eje $y$: $(0,-2)$.

\textbf{Paso 3.} Interpretamos $m=\frac12$: desde $(0,-2)$ avanzamos dos
unidades a la derecha y subimos una. Así obtenemos $(2,-1)$.

\textbf{Paso 4.} Trazamos la recta que pasa por ambos puntos.
\begin{center}
\begin{tikzpicture}
\begin{axis}[ga axis,width=8cm,height=5.7cm,xmin=-4,xmax=6,ymin=-5,ymax=3]
\addplot[domain=-4:6,samples=2,very thick,GArosa]{x/2-2};
\addplot[only marks,mark=*,GAazul] coordinates {(0,-2) (2,-1)};
\node[below left] at (axis cs:0,-2) {$(0,-2)$};
\node[above left] at (axis cs:2,-1) {$(2,-1)$};
\end{axis}
\end{tikzpicture}
\end{center}
\end{solution}
\end{example}

\subsection{Recta determinada por dos puntos}

Dos puntos distintos determinan una única recta. Primero se calcula la
pendiente y después se utiliza cualquiera de ellos en la forma
punto--pendiente.

\begin{formula}[title={Forma de dos puntos}]
\[
\boxed{y-y_1=\frac{y_2-y_1}{x_2-x_1}(x-x_1)},
\qquad x_1\neq x_2.
\]
\end{formula}

\begin{example}{Ecuación a partir de dos puntos}{dos-puntos}
Obtén la ecuación de la recta que pasa por $A(-1,4)$ y $B(3,-2)$.
\begin{solution}
\textbf{Paso 1.} Calculamos la pendiente:
\[
m=\frac{-2-4}{3-(-1)}=\frac{-6}{4}=-\frac32.
\]

\textbf{Paso 2.} Usamos el punto $A(-1,4)$:
\[
y-4=-\frac32(x+1).
\]

\textbf{Paso 3.} Despejamos $y$:
\[
y=-\frac32x+\frac52.
\]

\textbf{Paso 4.} Comprobamos con $B(3,-2)$:
$-\frac32(3)+\frac52=-2$.
\end{solution}
\end{example}

\subsection{Rectas horizontales y verticales}

\begin{properties}
\[
\boxed{y=k}\quad\text{es horizontal},
\qquad
\boxed{x=h}\quad\text{es vertical}.
\]
La primera pasa por todos los puntos $(x,k)$; la segunda, por todos los puntos
$(h,y)$. Una recta vertical no puede escribirse como $y=mx+b$ porque su
pendiente no está definida.
\end{properties}

\begin{exercise}{Construcción de ecuaciones}{ej-ecuaciones}
Obtén la ecuación de cada recta en forma pendiente--ordenada cuando sea
posible.
\begin{enumerate}[label=\alph*)]
  \item Pendiente $3$ y punto $(2,-1)$.
  \item Pendiente $-\frac12$ y punto $(-4,5)$.
  \item Pasa por $(1,2)$ y $(5,10)$.
  \item Pasa por $(-3,4)$ y $(2,4)$.
  \item Pasa por $(6,-2)$ y $(6,7)$.
\end{enumerate}
\end{exercise}

\section{Formas general y simétrica}

\subsection{Forma general}

\begin{definition}{Ecuación general de una recta}{forma-general}
La ecuación
\[
\boxed{Ax+By+C=0},
\]
con $A$ y $B$ no ambos nulos, representa una recta.
\end{definition}

Si $B\neq0$, puede despejarse $y$:
\[
y=-\frac ABx-\frac CB.
\]
Por tanto, su pendiente es $m=-\frac AB$ y su ordenada al origen es
$b=-\frac CB$.

\begin{example}{Conversión entre formas}{conversion}
Convierte $3x-2y-8=0$ a la forma pendiente--ordenada e identifica sus
parámetros.
\begin{solution}
\textbf{Paso 1.} Aislamos el término que contiene $y$:
\[
-2y=-3x+8.
\]

\textbf{Paso 2.} Dividimos entre $-2$:
\[
y=\frac32x-4.
\]

\textbf{Paso 3.} Leemos los parámetros: $m=\frac32$ y $b=-4$.
\end{solution}
\end{example}

\subsection{Intersecciones con los ejes y forma simétrica}

Para hallar el corte con el eje $x$ se hace $y=0$; para hallar el corte con el
eje $y$ se hace $x=0$.

Si la recta corta los ejes en $(a,0)$ y $(0,b)$, con $a,b\neq0$, puede
escribirse así:

\begin{formula}[title={Forma simétrica o de interceptos}]
\[
\boxed{\frac xa+\frac yb=1}.
\]
\end{formula}

\begin{example}{Interceptos y forma simétrica}{interceptos}
Escribe $2x+3y-12=0$ en forma simétrica.
\begin{solution}
\textbf{Paso 1.} Corte con $x$: hacemos $y=0$.
\[
2x-12=0\Rightarrow x=6.
\]

\textbf{Paso 2.} Corte con $y$: hacemos $x=0$.
\[
3y-12=0\Rightarrow y=4.
\]

\textbf{Paso 3.} Sustituimos $a=6$ y $b=4$:
\[
\boxed{\frac x6+\frac y4=1}.
\]
\end{solution}
\end{example}

\begin{exercise}{Formas de la recta}{ej-formas}
\begin{enumerate}[label=\alph*)]
  \item Convierte $y=-3x+7$ a forma general.
  \item Despeja $y$ en $4x+5y-20=0$ e identifica $m$ y $b$.
  \item Obtén los interceptos de $3x-2y-6=0$.
  \item Escribe en forma general la recta $\frac{x}{5}+\frac{y}{-2}=1$.
  \item Explica por qué $x=4$ no tiene forma pendiente--ordenada.
\end{enumerate}
\end{exercise}

\section{Paralelismo y perpendicularidad}

\begin{properties}
Sean dos rectas no verticales de pendientes $m_1$ y $m_2$:
\[
\boxed{m_1=m_2}\quad\Longrightarrow\quad\text{rectas paralelas},
\]
\[
\boxed{m_1m_2=-1}\quad\Longrightarrow\quad\text{rectas perpendiculares}.
\]
Las rectas verticales son paralelas entre sí y perpendiculares a las
horizontales.
\end{properties}

\begin{example}{Recta paralela por un punto}{paralela}
Obtén la recta que pasa por $P(2,-3)$ y es paralela a $3x-y+5=0$.
\begin{solution}
\textbf{Paso 1.} Despejamos $y$ en la recta dada:
\[
y=3x+5.
\]
Su pendiente es $m=3$.

\textbf{Paso 2.} La recta buscada debe tener la misma pendiente. Usamos
$P(2,-3)$:
\[
y+3=3(x-2).
\]

\textbf{Paso 3.} Simplificamos:
\[
\boxed{y=3x-9}.
\]
\end{solution}
\end{example}

\begin{example}{Recta perpendicular por un punto}{perpendicular}
Obtén la recta que pasa por $Q(-1,4)$ y es perpendicular a
$2x+5y-10=0$.
\begin{solution}
\textbf{Paso 1.} Identificamos la pendiente de la recta dada:
\[
m_1=-\frac AB=-\frac25.
\]

\textbf{Paso 2.} Buscamos su recíproco negativo:
\[
m_2=\frac52,qquad m_1m_2=-1.
\]

\textbf{Paso 3.} Usamos $Q(-1,4)$:
\[
y-4=\frac52(x+1).
\]
Esta es la ecuación solicitada; en forma general es
\[
\boxed{5x-2y+13=0}.
\]
\end{solution}
\end{example}

\begin{exercise}{Paralelas y perpendiculares}{ej-posicion}
\begin{enumerate}[label=\alph*)]
  \item Decide si $y=4x-1$ y $8x-2y+7=0$ son paralelas.
  \item Decide si $3x+y=2$ y $x-3y=5$ son perpendiculares.
  \item Obtén la paralela a $2x-3y+6=0$ que pasa por $(3,1)$.
  \item Obtén la perpendicular a $y=-\frac14x+2$ que pasa por $(-2,5)$.
  \item Encuentra el valor de $k$ para que $kx+2y-1=0$ sea paralela a
  $3x-y+4=0$.
\end{enumerate}
\end{exercise}

\section{Intersección de dos rectas}

El punto de intersección de dos rectas satisface simultáneamente ambas
ecuaciones. Por ello se obtiene resolviendo un sistema lineal de dos
ecuaciones con dos incógnitas.

\begin{example}{Intersección mediante sustitución}{interseccion}
Encuentra la intersección de
\[
r_1:y=2x-1,
\qquad
r_2:x+y=8.
\]
\begin{solution}
\textbf{Paso 1.} Sustituimos $y=2x-1$ en la segunda ecuación:
\[
x+(2x-1)=8.
\]

\textbf{Paso 2.} Resolvemos para $x$:
\[
3x=9\Rightarrow x=3.
\]

\textbf{Paso 3.} Sustituimos en cualquiera de las rectas:
\[
y=2(3)-1=5.
\]

\textbf{Paso 4.} Comprobamos: $3+5=8$. La intersección es
$\boxed{(3,5)}$.
\end{solution}
\end{example}

\begin{remark}
Dos rectas con pendientes diferentes tienen una intersección. Dos rectas
paralelas distintas no tienen solución común; si sus ecuaciones representan
la misma recta, tienen infinitas soluciones.
\end{remark}

\begin{exercise}{Intersecciones}{ej-interseccion}
Encuentra el punto de intersección o describe el tipo de sistema.
\begin{enumerate}[label=\alph*)]
  \item $y=3x-2$ y $y=-x+6$.
  \item $2x+y=7$ y $x-y=2$.
  \item $4x-2y=8$ y $2x-y=1$.
  \item $3x+6y=9$ y $x+2y=3$.
\end{enumerate}
\end{exercise}

\section{Ángulo entre dos rectas}

El ángulo menor $\varphi$ entre dos rectas se obtiene a partir de sus
pendientes. La expresión proviene de la identidad de la tangente de una
diferencia de ángulos.

\begin{formula}[title={Ángulo entre dos rectas}]
\[
\boxed{
\tan\varphi=\left|\frac{m_2-m_1}{1+m_1m_2}\right|
},
\qquad 0^\circ\leq\varphi\leq90^\circ.
\]
Si $1+m_1m_2=0$, las rectas son perpendiculares y
$\varphi=90^\circ$.
\end{formula}

\begin{example}{Cálculo del ángulo menor}{angulo-rectas}
Calcula el ángulo entre $r_1:y=2x+1$ y $r_2:y=-\frac13x+4$.
\begin{solution}
\textbf{Paso 1.} Identificamos $m_1=2$ y $m_2=-\frac13$.

\textbf{Paso 2.} Sustituimos:
\[
\tan\varphi=\left|
\frac{-\frac13-2}{1+2(-\frac13)}
\right|
=\left|\frac{-\frac73}{\frac13}\right|=7.
\]

\textbf{Paso 3.} Aplicamos la función inversa:
\[
\varphi=\arctan(7)\approx81.87^\circ.
\]
\end{solution}
\end{example}

\begin{exercise}{Ángulos}{ej-angulos}
Calcula el ángulo menor entre cada pareja.
\begin{enumerate}[label=\alph*)]
  \item $y=x+1$ y $y=-x+3$.
  \item $y=\frac12x-4$ y $y=2x+1$.
  \item $3x-y+2=0$ y $x+2y-5=0$.
\end{enumerate}
\end{exercise}

\section{Distancia de un punto a una recta}

La distancia de un punto a una recta es la longitud del segmento
perpendicular que los une. No es la distancia a cualquier punto elegido sobre
la recta.

\begin{formula}[title={Distancia de un punto a una recta}]
Si $P(x_0,y_0)$ y la recta es $Ax+By+C=0$, entonces
\[
\boxed{
d(P,r)=\frac{|Ax_0+By_0+C|}{\sqrt{A^2+B^2}}
}.
\]
\end{formula}

\begin{example}{Distancia punto--recta}{distancia-recta}
Calcula la distancia de $P(4,-1)$ a la recta $3x-4y-6=0$.
\begin{solution}
\textbf{Paso 1.} Identificamos
$A=3$, $B=-4$, $C=-6$, $x_0=4$ y $y_0=-1$.

\textbf{Paso 2.} Sustituimos en el numerador:
\[
|3(4)-4(-1)-6|=|12+4-6|=10.
\]

\textbf{Paso 3.} Calculamos el denominador:
\[
\sqrt{3^2+(-4)^2}=\sqrt{25}=5.
\]

\textbf{Paso 4.} Dividimos:
\[
d(P,r)=\frac{10}{5}=\boxed{2}.
\]
\end{solution}
\end{example}

\begin{exercise}{Distancia punto--recta}{ej-distancia-recta}
\begin{enumerate}[label=\alph*)]
  \item Distancia de $(2,3)$ a $4x-3y+5=0$.
  \item Distancia del origen a $5x+12y-26=0$.
  \item Encuentra los valores de $k$ para que $(k,1)$ esté a distancia $2$
  de la recta $x=4$.
  \item Verifica que un punto que satisface $Ax+By+C=0$ tiene distancia cero
  a la recta.
\end{enumerate}
\end{exercise}

\section{Problemas de unidad}

\begin{problem}{Aplicaciones de la línea recta}{problemas-recta}
\begin{enumerate}
  \item La temperatura de un líquido aumenta linealmente. A los $2$ minutos
  es $18\,{}^\circ\mathrm C$ y a los $8$ minutos es
  $42\,{}^\circ\mathrm C$. Obtén el modelo $T(t)$ e interpreta su pendiente.
  \item Una carretera recta pasa por los puntos cartográficos $A(-3,2)$ y
  $B(5,6)$. Obtén su ecuación general.
  \item Determina la ecuación de la mediatriz del segmento cuyos extremos son
  $P(-2,1)$ y $Q(6,5)$.
  \item Los lados de un triángulo están sobre $x+y=6$, $x=1$ y $y=2$.
  Obtén sus vértices y su área.
  \item Encuentra la distancia entre las rectas paralelas
  $3x-4y+5=0$ y $3x-4y-15=0$.
  \item Una recta pasa por $(2,-1)$ y forma un ángulo de $45^\circ$ con el
  eje $x$. Obtén su ecuación.
  \item Encuentra $k$ para que las rectas $kx+3y-1=0$ y
  $2x-ky+4=0$ sean perpendiculares.
\end{enumerate}
\end{problem}

\begin{hints}
Para una mediatriz, encuentra primero el punto medio y después una pendiente
perpendicular. Para la distancia entre rectas paralelas, toma un punto de una
de ellas y calcula su distancia a la otra.
\end{hints}

\clearpage
\section*{Resumen}
\addcontentsline{toc}{section}{Resumen}

\begin{properties}
\[
\begin{aligned}
m&=\frac{y_2-y_1}{x_2-x_1},
&m&=\tan\theta,\\[1mm]
y-y_1&=m(x-x_1),
&y&=mx+b,\\[1mm]
Ax+By+C&=0,
&\frac xa+\frac yb&=1,\\[1mm]
m_1=m_2&\Longrightarrow r_1\parallel r_2,
&m_1m_2=-1&\Longrightarrow r_1\perp r_2.
\end{aligned}
\]
\[
\tan\varphi=\left|\frac{m_2-m_1}{1+m_1m_2}\right|,
\qquad
d(P,r)=\frac{|Ax_0+By_0+C|}{\sqrt{A^2+B^2}}.
\]
\end{properties}

\section*{Comprobación rápida}
\addcontentsline{toc}{section}{Comprobación rápida}

\begin{enumerate}
  \item Calcula la pendiente entre $(-2,3)$ y $(4,-6)$.
  \item Obtén la recta de pendiente $2$ que pasa por $(1,-4)$.
  \item Convierte $2x-5y+10=0$ a forma pendiente--ordenada.
  \item Obtén la perpendicular a $y=-3x+1$ que pasa por el origen.
  \item Encuentra la intersección de $y=x+2$ y $y=-2x+8$.
  \item Calcula la distancia de $(1,1)$ a $3x+4y-12=0$.
\end{enumerate}

\begin{hints}[title={Resultados}]
$-\frac32$; \quad $y=2x-6$; \quad $y=\frac25x+2$; \quad
$y=\frac13x$; \quad $(2,4)$; \quad $1$.
\end{hints}
